\section{Instinct, Consciousness, Science}

Corresponding to the three stages of the consciousness of the individual as described in Julius Evola's \textit{The Individual and the Becoming of the World}, there are three stages in the self-understanding of the social group:

\begin{itemize}


\item Instinct

\item Consciousness

\item Science
\end{itemize}


\paragraph{Instinct}


A healthy society is based on healthy instincts. Families are strong, children are valued, elders are respected, ancestors are glorified, economic activity leads to near self-sufficiency, folk religion creates unity, and threats --- whether internal or external --- are resisted without question. This is clear from the fundamental tenet of Hermetism: ``as above, so below''. Since manifestation on the physical plane is the reflection of the spirit, a healthy society follows from a healthy spirit. Each race or national group is connected with its own peculiar spirit. At the level of instinct, this connection is implicit and does not arise to the level of consciousness.

\paragraph{Consciousness}


For some members of the society, the values and traditions begin to come into consciousness and may be questioned. In a healthy society, there a system of religious dogma and science that will aid these members in their self-understanding. For example, in \textbf{Campanella}'s City of the Sun\footnote{\url{https://gornahoor.net/?p=518}}, there is training in all the skills and values required for the well-ordering of the city. Then there is training in the various sciences and metaphysics for those with the interest ability. This is the level of Traditional sciences and exoteric religion as described by \textbf{Rene Guenon}. So, while the first group is grounded on instinct, this group is grounded on science and religion. In a Traditional society, there is no opposition between the two groups; the second simply has a deeper and more explicit self-understanding.

In the hands of revolutionaries, so-called ``consciousness raising'' has the opposite effect. Instead of reinforcing Traditional wisdom, it seeks to undermine the natural structure of society. In the hands of a \textbf{Friedrich Nietzsche}, the ``hermeneutics of suspicion'' is effective in exposing the hypocrisies of an unhealthy society. However, the revolutionary seeks to undermine the healthy aspects of society. As tools, he encourages pornography, sexual immorality, hatred for the faults of the ancestors, disrespect for legitimate authority, class resentment, the devaluation of marriage and family, egalitarianism, and ridicule of spiritual reality and religion.

Unfortunately, we often find those who call themselves ``Traditionalists'' or claim to be on the ``right'', expounding the same values as the revolutionaries. This merely justifies Julius Evola's claim that revolutionary ideas have penetrated down to the roots of the mind of Western man, to the extent they don't realize it and cannot find any way out.

\paragraph{Science}


By ``Science'' we mean an organized body of knowledge. We may call those at the second level the consumers of science, and those at this level its creators. Beyond the practical sciences required for the physical and economic well-being of the group, there are the Traditional sciences and, of course, Metaphysics proper, which is as precise and exact as any science. Those in this group are fully conscious of the sources of the values, mores, and existence of the society. Ideally, they rule the society. Through education and religion, they maintain the health of the whole.

Note that there is no class warfare in a healthy society. The elite rule for the benefit of all; in return, they benefit from the free allegiance of the workers, and can thus pursue their scientific and metaphysical interests. I have heard ``new rightists'' speak of the masses in very uncomplimentary terms. It is true that the masses are malleable; that is what makes the revolutionary so dangerous. But the elite are supposed to be the consciousness of the masses, not its enemy.

By understanding this schema of Instinct, Consciousness, and Science, all social phenomena can be understood. There are three functions of the mind: Thinking, Willing, Feeling. At the level of Consciousness, the mind is restricted to thinking. That is why it is so sterile and often unrelated to reality. This makes it prone to the absorption of ideas, almost willy-nilly, and the lack of a True Will and strong feeling makes them hollow and shallow. Unfortunately, in the modern West, the dominant intellectuality is at this level. Only at the level of Instinct and Science is the full range of mind in play. What the common man knows in his gut, the true elite knows in his heart. The revolutionary is aware of this, and makes both his enemy.


\flrightit{Posted on 2010-06-10 by Cologero}

